\documentclass[11pt]{amsart}

\usepackage[utf8]{inputenc}
\usepackage[T1]{fontenc}
\usepackage{amsmath,amssymb}
\usepackage{booktabs}
\usepackage{url}

\title{Ablation study, experiment 1: isolating the CA framing}
\author{Kenza Benjelloun}
\date{28 August 2026}

\begin{document}

\maketitle

\section{Position in the chapter}

Experiment 1 is the first ablation of the final chapter. It isolates a single
factor: whether the cellular automaton framing --- feeding each cell its
spatial neighborhood --- contributes anything beyond the temporal recurrence a
plain RNN already provides.

Everything else is held fixed. Both \emph{streams} use the same architecture, the same
sequence length, the same splits, the same normalisation statistics and the
same underlying grid. The only difference is the shape of the input tensor:

\begin{itemize}
  \item \texttt{nn} --- $(N, 6, 1)$: each sample is one cell's own history
    over six days. Recurrence only.
  \item \texttt{nca} --- $(N, 6, 9)$: each sample is the same six days, but
    each timestep carries the cell's $3 \times 3$ Moore neighborhood.
    Recurrence plus CA-style spatial context.
\end{itemize}

The ablation is therefore of the second axis of the tensor's last dimension:
$F = 1$ against $F = 9$. If \texttt{nca} does not beat \texttt{nn}, the CA
framing is not earning the ninefold increase in input width, and the chapter's
central claim needs qualifying.

The data for both streams is generated and verified. This document records how,
and the one problem that had to be fixed first.

\section{Generation}

Produced by \texttt{src/data/prep.py}, which shares all loading, masking,
splitting and saving logic between the two streams and differs only in the
extractor it calls. That sharing is deliberate: it is what allows any
difference in the results to be attributed to the framing rather than to
preprocessing.

Dataset \texttt{cmems\_mod\_med\_phy-temp\_my\_4.2km\_P1D-m}, version 202511,
variable \texttt{thetao} (SST, ${}^\circ$C). Bounding box longitude 12--16,
latitude 44.5--45.5 (Northern Adriatic), surface level only. Sequence length 6.

\begin{table}[h]
\begin{tabular}{llrr}
\toprule
Split & Range & Mean & Std \\
\midrule
train & 2021-01-01 $\to$ 2023-12-31 & 17.8304 & 5.8792 \\
val   & 2024-01-01 $\to$ 2024-12-31 & 18.6804 & 5.8236 \\
test  & 2025-01-01 $\to$ 2026-07-31 & 17.6932 & 5.9215 \\
\bottomrule
\end{tabular}
\end{table}

Disjoint and chronologically ordered. The test window ends at the last day of
product coverage; the final six days yield no target, so the effective test
window ends around 2026-07-25. The train statistics are to be used for all
three splits in both streams --- normalising val or test by their own statistics
would leak distributional information and would do so asymmetrically between
streams of the experiment.

These statistics vary in the fourth decimal place between machines (demetra
reported 17.8332 / 5.8655 for train). The cause is floating-point accumulation
order in the dask reduction, not a difference in the data. The value embedded
in \texttt{1\_nca\_train.pt} is the one to use, so that a single number is
applied consistently rather than recomputed per run.

\begin{table}[h]
\begin{tabular}{lrrr}
\toprule
File & $X$ shape & $Y$ shape & Size \\
\midrule
\texttt{1\_nn\_train.pt}  & $(806949, 6, 1)$ & $(806949, 1)$ & 22.6 MB \\
\texttt{1\_nn\_val.pt}    & $(266760, 6, 1)$ & $(266760, 1)$ & 7.5 MB \\
\texttt{1\_nn\_test.pt}   & $(423111, 6, 1)$ & $(423111, 1)$ & 11.8 MB \\
\texttt{1\_nca\_train.pt} & $(806949, 6, 9)$ & $(806949, 1)$ & 177.5 MB \\
\texttt{1\_nca\_val.pt}   & $(266760, 6, 9)$ & $(266760, 1)$ & 58.7 MB \\
\texttt{1\_nca\_test.pt}  & $(423111, 6, 9)$ & $(423111, 1)$ & 93.1 MB \\
\bottomrule
\end{tabular}
\end{table}

All float32, zero NaNs, value ranges 5.4--31.5${}^\circ$C. Each file carries a
\texttt{config} dictionary alongside $X$ and $Y$ recording the experiment
number, feature variant, split, date range, dataset id, bounding box and
normalisation statistics, so that a file can be identified independently of
its name.

\section{The control was broken, and is now fixed}

As first generated, the two streams were not on the same sample set. The
\texttt{nca} stream retained $79.42\%$ of the samples the \texttt{nn} stream
did, identically in every split: 1016037 against 806949 for train, 335880
against 266760 for val, 532743 against 423111 for test.

Both extractors crop to the same interior cells, so the candidate sets were
identical. The divergence came from the NaN mask, which dropped any sample with
a missing value anywhere in its features. For \texttt{nn} that condition
involved the centre cell alone; for \texttt{nca} it involved all nine. Every
cell with a land or missing neighbour therefore survived in \texttt{nn} and was
dropped in \texttt{nca}: the whole coastline ring.

This defeated the purpose of the ablation. The dropped $20.58\%$ was not a
random subset but precisely the coast-adjacent cells, which are plausibly the
hardest to predict --- shallower water, stronger diurnal and terrestrial
forcing. The \texttt{nn} stream was thus scored on a harder cell population than
the \texttt{nca} stream, and any MAE gap between them would have confounded the
effect of the CA framing with the effect of a different evaluation set. The
confound ran in the direction that flatters \texttt{nca}, which is the direction
the chapter's claim would benefit from, so it could not be left unaddressed.

That the ratio was identical to four decimal places across three independent
splits confirmed the cause was fixed geometry, a land mask, and not anything
data-dependent.

\subsection*{Options considered}
\begin{enumerate}
  \item Apply the \texttt{nca} validity mask to both experiment streams. Restricts
    \texttt{nn} to the same 806949 cells, restores exact pairing and makes
    per-cell paired tests possible. Costs $20\%$ of the \texttt{nn} data,
    which is not a real loss here since the comparison is the point.
  \item Keep both as generated, report the cell counts explicitly and refuse
    any cross-stream comparison. This effectively abandons experiment 1.
  \item Pad or impute missing neighbours so that \texttt{nca} retains the
    coastal cells. Relates to the existing TODO in \texttt{extractors.py} on
    non-full neighborhoods, and raises its own question of what a synthetic
    neighbour means physically at a coastline.
  \item Think deeper and come up with a strategy to deal better with land cells.
  This was a weak point which was also pointed out by a referee on the paper
  published in \(\mathrm{ACRI}\).
\end{enumerate}

\subsection*{Resolution}

Option 1 is implemented. The mask is now computed by a single function,
\texttt{neighborhood\_valid\_mask}, which derives validity from the grid rather
than from the feature tensor: a cell is valid if and only if its full
$3 \times 3$ neighborhood is finite at every timestep and its target is finite.
Both extractors receive the same boolean mask. The criterion is deliberately
stricter than \texttt{nn} requires, since \texttt{nn} only ever reads the centre
cell.

The point of computing the mask from the grid rather than from $X_t$ is that
masking on $X_t$ makes the criterion depend on the feature count, which is
exactly what produced the divergence. The function carries a comment recording
this, so that the apparently redundant second \texttt{unfold} is not optimised
away later.

Sample counts now match exactly: 806949, 266760 and 423111 for both streams,
in the same row-major order, so index $i$ refers to the same grid point in both
tensors and per-cell paired tests are available.

Two secondary observations corroborate the fix. The \texttt{nn} value ranges now
sit strictly inside the \texttt{nca} ranges --- train $X$ minimum rises from
5.388 to 5.814 --- because the extreme values previously present in \texttt{nn}
came from coastal cells that no longer survive. And \texttt{nn}'s $X$ range now
coincides exactly with its $Y$ range in every split, as it must when both are
drawn from the same cell set.

Options 3 and 4 remain open, and belong in a later experiment: they change what
\texttt{nca} denotes rather than holding it fixed. Option 4 is the most
interesting one, but I doubt I'll have time to think about it. You're welcome to
contribute any ideas you might have.

\section{Pipeline changes}

The download is now separated from the extraction. \texttt{src/data/download.py}
fetches each split once to local netCDF under \texttt{data/raw/} and skips files
that already exist; \texttt{--force} re-fetches. \texttt{prep.py} reads from
those local files and performs no network access. Running both feature variants
therefore costs one download rather than two, and re-running an extraction after
a code change costs none at all.

Inspecting the downloaded files revealed that the subset was being fetched with
all 141 depth levels, of which only the surface is ever used --- 1.4 GB per split
where tens of megabytes suffice. The download now bounds depth explicitly. The
selected level is unchanged, and the regenerated tensors are identical in shape
and sample count to those produced from the full-depth files.

The pipeline also runs locally now, in a conda environment (\texttt{nca},
Python 3.11, CPU-only torch) specified by \texttt{env/carnn\_env\_min.yml}. A
full transitive pin is recorded alongside it in \texttt{env/carnn\_env.yml}.
Generation takes a few minutes per stream on the laptop, which removes the
dependency on cluster queue availability for iteration. Training runs still
belong on demetra.

To reproduce: \texttt{python -m src.data.download}, then
\texttt{python -m src.data.prep --experiment 1 --features nca} and again with
\texttt{--features nn}. On demetra, \texttt{sbatch src/data/prep.slurm} from the
repository root with \texttt{carnn\_venv} active.

\end{document}